\section{Connectedness}

One interesting motivating example for studying connectedness is the following result from calculus:

\begin{theorem}[Intermediate Value Theorem]
    Let $f:[0,1]\to\R$ be a continuous function. WLOG, we assume that $f(0)<f(1)$. Then, for any $r\in [f(0),f(1)]$, there exists $c\in (0,1)$ such that $f(c)=r$.
\end{theorem}

This is a consequence of the fact that $[0,1]$ is ``connected''.

% TODO: include examples and counterexamples, discuss the components of a
% space and the relationship between connectedness and path‑connectedness.

%---------------------------------------------------------------------------
% Section 9: Homotopy and the fundamental group (Willard §§32–34)
%---------------------------------------------------------------------------

\subsection{Connected spaces}


\begin{definition}
Let $X$ be a topological space. A \emph{separation} of $X$ is a pair of disjoint nonempty open sets $U,V\subseteq X$ with $X=U\cup V$. The space $X$ is \emph{disconnected} if it admits a separation; otherwise $X$ is \emph{connected}.
\end{definition}

\begin{remark}
    This is a pure topological property: it depends only on the open sets of $X$. In particular, connectedness is preserved by homeomorphisms.
\end{remark}

\begin{remark}
    $X$ is connected iff the only subsets that are both open and closed (clopen) are $\varnothing$ and $X$.
\end{remark}



\begin{example}
    $[-1,0)\cup (0,1]$ is disconnected.
\end{example}

\begin{example}
    $\{a,b\}$ with the trivial topology $\{\emptyset,\{a,b\}\}$ is connected.
\end{example}

\begin{example}
    $\Q$ is disconnected: for $r$ irrational, $(-\infty,r)\cap \Q$ and $(r,\infty)\cap \Q$ separate $\Q$.
\end{example}


How do we check if a space is connected? 

\begin{theorem}
    If $f:X\to Y$ is continuous and $X$ is connected, then $f(X)$ is connected.
\end{theorem}
\begin{proof}
    If $U,V$ separate $f(X)$, then $f^{-1}(U), f^{-1}(V)$ separate $X$.
\end{proof}

\begin{lemma}
    If $X$ has separation $U, V$ and $Y\subseteq X$ is connected, then $Y\subset U$ or $Y\subset V$.
\end{lemma}
\begin{proof}
    If not, then $Y\cap U$ and $Y\cap V$ are nonempty, disjoint, open in the subspace topology on $Y$, and cover $Y$, contradicting connectedness of $Y$.
\end{proof}

\begin{theorem}
    If $A\subset X$ is connected and $A\subseteq B\subseteq \overline{A}$, then $B$ is connected.
\end{theorem}

\begin{proof}
    Suppose $B$ has a separation $C,D$. Then, by the lemma, $A\subseteq C$ or $A\subseteq D$. WLOG, $A\subseteq C$. Then, $\overline{A}\subseteq \overline{C}$. Note that $\overline{C}_B=B\cap \overline{C}$. Since $C$ is closed in $B$, we have $\overline{C}_B=C$. Thus, $\overline{C}\cap B=C$. Hence, $\overline{A}\cap B\subseteq C$, so $B\subseteq C$, contradicting that $D$ is nonempty. 
\end{proof}

\begin{theorem}
    The union of a collection of connected subspaces of $X$ that have a point in common is connected.
\end{theorem}
\begin{proof}
    Let $Y=\bigcup_\alpha Y_\alpha$ with $y\in Y_\alpha\subset X$ for all $\alpha$. 

    Assume that $Y$ has a separation $U,V$. Suppose $y\in U$. Then, by the lemma, $Y_\alpha\subseteq U$ for all $\alpha$. Thus, $Y\subseteq U$, contradicting that $V$ is nonempty.
\end{proof}

\begin{theorem}
    A finite product of connected spaces is connected.
\end{theorem}
\begin{proof}
    We prove the case of two spaces $X,Y$. Let $(a,b)\in X\times Y$. For $x\in X$, let $T_x=(X\times \{b\})\cup (\{x\}\times Y)$. Note that $T_x$ is connected by the previous theorem and the fact that the product of a space with a singleton is homeomorphic to the space itself.

    Note that $X\times Y=\cup_{x\in X}T_x$ is connected again by the previous theorem, since all $T_x$ contain the point $(a,b)$.

    The general case follows by induction.
\end{proof}

\begin{example}
    $\R^n$ is connected for all $n\geq 1$ (if we assume the connectedness of $\R$ which will be proved later).
\end{example}

\begin{example}
    $\R^\N$ with box topology is not connected. Indeed, let
    \[U=\{\text{all bounded sequences}\},\quad V=\{\text{all unbounded sequences}\}.\]
    Then, $U$ and $V$ are disjoint, nonempty and $U\cup V=\R^\N$. Given $\mathbf{a}\in \R^\N$, note that $(a_1-1,a_1+1)\times (a_2-1,a_2+1)\times\cdots$ is an open neighborhood of $\mathbf{a}$ contained in $U$ (resp. $V$) is $\mathbf{a}$ is bounded (resp. unbounded). Hence,  $U$ and $V$ are open in the box topology, so they separate $\R^\N$.
\end{example}

\begin{example}
    $\R^\N$ with product topology is connected. 
    
    Let $\tilde{\R}^n=\{x\in\R^\N:x_i=0\,\forall i>n\}$.
    Note that $\tilde{\R}^n\cong \R^n$ so $\tilde{\R}^n$ is connected. Let $\R^\infty=\cup_{n\geq 1}\tilde{\R}^n$. Then, $\R^\infty$ is connected as the union of connected sets with a point in common (the zero sequence). Finally, we claim that $\overline{\R^\infty}=\R^\N$ to conclude the proof. Indeed, for any $x\in\R^\N$ and any basic open neighborhood $U=\prod_{i\geq 1}U_i$ of $x$ (where $U_i=\R$ for all $i>N$ for some $N$), we have that $(x_1,x_2,\ldots,x_N,0,0,\cdots)\in \R^\infty\cap U$. Hence, $x\in\overline{\R^\infty}$.
\end{example}

The same strategy can be used to show that arbitrary products of connected spaces are connected in the product topology.


\subsection{Connected subspaces of $\R$}

We shall prove that $\R$ is connected and so are intervals and rays in $\R$. It turns out that this only depends on the order structure of $\R$.

\begin{definition}
    A linear poset $L$ having more than one element is called a linear continuum if the following hold:
    \begin{enumerate}
        \item $L$ has the least upper bound property: every nonempty subset of $L$ that is bounded above has a least upper bound. (sup exists)
        \item If $x<y$, then there exists  $z\in L$ such that $x<z<y$.
    \end{enumerate}
\end{definition}

\begin{definition}
    $Y\subset L$ is called convex if for any $a,b\in Y$ with $a<b$, we have that $[a,b]\subseteq Y$.
\end{definition}

\begin{theorem}
    If $L$ is a linear continuum with the order topology, then every convex subset $Y$ of $L$ is connected.
\end{theorem}
\begin{proof}
    Assume that $Y$ has a separation $A,B$. Pick $a\in A$ and $b\in B$ with $a<b$. Then, $[a,b]\subset Y$. Let $A_0=A\cap [a,b]$ and $B_0=B\cap [a,b]$. Then, $A_0,B_0$ are open in $[a,b]$ w.r.t. the subspace topology which is the same as the order topology on $[a,b]$. So $A_0,B_0$ form a separation of $[a,b]$.

    Let $c=\sup A_0$. We show that $c\notin A_0$ and $c\notin B_0$, contradicting that $[a,b]=A_0\cup B_0$.

    Suppose $c\in B_0$. Then, $c\neq a$, so $c=b$ or $a<c<b$. Since $B_0$ is open, in either case, there exists $d$ such that $(d,c]\subset B_0$. If $c=b$, then $d$ is an upper bound of $A_0$ less than $c$, contradicting that $c=\sup A_0$. If $a<c<b$, then since $c$ is an upper bound, $(c,b]\cap A_0=\emptyset.$ Hence, $(d,b]=(d,c]\cup (c,b]\cap A_0=\emptyset$ which makes $d$ an upper bound of $A_0$ less than $c$, again a contradiction. Thus, $c\notin B_0$.

    Now suppose $c\in A_0$. Then, $c\neq b$, and so $c=a$ or $a<c<b$. Since $A_0$ is open, in either case, there exists $e$ such that $[c,e)\subset A_0$. But then any $z$ such that $c<z<e$ is in $A_0$, contradicting that $c=\sup A_0$ ($z$ exists because of (2)). Thus, $c\notin A_0$.
\end{proof}


Now we can prove the intermediate value theorem.
\begin{theorem}[Intermediate Value Theorem]
    Let $f:[a,b]\to\R$ be a continuous function with $X$ being connected. Assume that $f(a)<f(b)$. Then, for any $r\in (f(a),f(b))$, there exists $c\in (a,b)$ such that $f(c)=r$.
\end{theorem}
\begin{proof}
    Note that $f([a,b])$ is connected since $[a,b]$ is connected. We now show that  $[f(a),f(b)]\subseteq f([a,b])$. Suppose not, then there exists $r\in (f(a),f(b))$ with $r\notin f([a,b])$. Let $U=(-\infty,r)$ and $V=(r,\infty)$. Then, $U\cap f([a,b])$ and $V\cap f([a,b])$ separate $f([a,b])$, contradicting connectedness of $f([a,b])$.
\end{proof}

The interval $[a,b]$ can be replaced by any connected space $X$.
\begin{theorem}[Intermediate Value Theorem]
    Let $f:X\to\R$ be a continuous function with $X$ being connected. Assume that $a,b\in X$ and that $f(a)<f(b)$. Then, for any $r\in (f(a),f(b))$, there exists $c\in X$ such that $f(c)=r$.
\end{theorem}

That intervals in $\R$ are connected gives a sufficient condition for showing that a space is connected.



\begin{definition}
    A space $X$ is path connected if every $x,y\in X$ can be joined by a path, i.e., there exists a continuous $f:[a,b]\to X$ with $f(a)=x$ and $f(b)=y$.
\end{definition}

\begin{lemma}
    A path connected space $X$ is connected.
\end{lemma}

\begin{proof}
    Suppose $U,V$ separates $X$. Let $f:[a,b]\to X$ be a path from $x\in U$ to $y\in V$. Since $[a,b]$ is connected, $f([a,b])$ is connected. So by a theorem we proved previously, $f([a,b])$ is contained in $U$ or $V$, contradicting that $f(a)=x\in U$ and $f(b)=y\in V$.
\end{proof}

\begin{example}
    The ball $\mathbb{B}^n=\{x\in\R^n|\|x\|\leq 1\}$ is connected.

    The sphere $\mathbb{S}^{n-1}=\{x\in\R^n|\|x\|= 1\}$ is connected for $n\geq 2$.
\end{example}

\begin{example}
    The topologist's sine curve is the closure of the following set in $\R^2$:
    \[
        S=\left\{\left(x,\sin\frac{1}{x}\right)|x\in (0,1]\right\}
    \]
    Note that
    \[
        \overline{S}=S\cup \{(0,y)|y\in [-1,1]\}.
    \]
    This space is connected but not path connected.

    To see that it is connected, note that $S$ is path connected (hence connected) and hence $\overline{S}$ is connected.

    To see that it is not path connected, suppose there exists a path $f:[a,b]\to \overline{S}$ with $f(a)=(0,0)$ and $f(b)=(1,\sin 1)$. Consider $A=f^{-1}(0\times [-1,1])\subset [a,b]$ which is the preimage of a closed set, hence closed. Let $c=\sup A\in A$. $c\neq b$ since $f(b)=(1,\sin 1)$. Then, $f((c,b])\subset S$. Replace $[c,b]$ by $[0,1]$ for convenience. We write $f(t)=(x(t),y(t))$. Then, $x(0)=0$ and $x(t)>0$ for all $t\in (0,1]$. Since $x$ is continuous, there exists $t_n\to 0$ with $x(t_n)=\frac{1}{n\pi/2}$ for all $n\geq 1$ (intermediate value theorem). Then, $y(t_n)=\sin(n\pi/2)=(-1)^n$ does not converge as $n\to\infty$, contradicting the continuity of $y$ at $0$.
\end{example}

% \begin{proposition}
% For a space $X$ the following are equivalent.
% \begin{enumerate}
%     \item $X$ is connected.
%     \item The only subsets of $X$ that are both open and closed (clopen) are $\varnothing$ and $X$.
%     \item Every continuous map $f:X\to\{0,1\}$ (discrete topology) is constant.
% \end{enumerate}
% \end{proposition}

% \begin{proof}
% (1)$\Rightarrow$(2): A nontrivial clopen set $H$ gives a separation $(H,X\setminus H)$.  
% (2)$\Rightarrow$(3): $\partial A=\varnothing$ implies $A$ is clopen.  
% (3)$\Rightarrow$(4): $f^{-1}(0)$ has empty boundary (preimages of points in a discrete codomain are clopen).  
% (4)$\Rightarrow$(1): If $X=H\cup K$ is a separation, the characteristic map $X\to\{0,1\}$ is nonconstant.  
% \end{proof}


% Connectedness is preserved by continuous images. This generalizes the intermediate value theorem.
% \begin{theorem}
% If $f:X\to Y$ is continuous and $X$ is connected, then $f(X)$ is connected. 
% \end{theorem}



\subsection{Components and local connectedness}

\begin{definition}
    Let $X$ be a topological space. We define an equivalence relation on $X$ by declaring $x\sim y$ iff there exists a connected subspace of $X$ containing both $x$ and $y$. The equivalence classes of this relation are called the \emph{components} of $X$.

    We define another equivalence relation on $X$ by declaring $x\approx y$ iff there exists a path in $X$ connecting $x$ and $y$. The equivalence classes of this relation are called the \emph{path components} of $X$.
\end{definition}

Symmetry and relexivity are clear. For transitivity, if $x,y$ lie in a connected subspace $A$ and $y,z$ lie in a connected subspace $B$, then $x,z$ lie in the connected subspace $A\cup B$ (since $y\in A\cap B$). 



\begin{theorem}
    The components of $X$ are connected disjoint subspaces of $X$ whose union is $X$, such that each nonempty connected subspace of $X$ intersects only one of them.
\end{theorem}

\begin{proof}
    Each connected subspace $A$ intersects only one of them: if $A$ intersects two components $C_1,C_2$, then pick $x_i\in A\cap C_i$ for $i=1,2$. Then, $x_1,x_2$ are equivalent, so $C_1=C_2$.

    To show that each component $C$ is connected, choose $x_0\in C$. For each $x\in C$, there exists a connected subspace $A_x$ containing $x_0$ and $x$. So $A_x\subset C$. Therefore, $C=\bigcup_{x\in C}A_x$ is connected as the union of connected sets with a point $x_0$ in common.
\end{proof}

\begin{theorem}
    The path components of $X$ are path connected disjoint subspaces of $X$ whose union is $X$, such that each nonempty path connected subspace of $X$ intersects only one of them.
\end{theorem}
The proof is similar to that of the previous theorem.

\begin{remark}
    Each component of $X$ is closed as its closure is connected. If $X$ has only finitely many components, then each component is also open, since its complement is a finite union of closed sets. This doesn't hold in general for infinitely many components.
\end{remark}

\begin{example}
    Each component of $\Q$ is a singleton which is not open. 
\end{example}

\begin{example}
    For the topologist's sine curve $\overline{S}$, the components are $\overline{S}$ itself, while the path components are $S$ and $V=\{(0,y)|y\in [-1,1]\}$. Note that $S$ is open but not closed in $\overline{S}$, while $V$ is closed but not open in $\overline{S}$.
\end{example}

Each path component is connected and hence contained in a component. When do they coincide?

\begin{definition}
    A space $X$ is \emph{locally (path) connected} iff for every $x\in X$ and every open neighborhood $U$ of $x$, there exists a (path) connected open neighborhood $V$ of $x$ with $V\subseteq U$.
\end{definition}

\begin{remark}
    locally path connected $\implies$ locally connected.
\end{remark}

\begin{example}
    Every open subset of $\R^n$ is locally path connected.
\end{example}

\begin{example}
    Let $X=\N$ with the complement finite topology. Then, $X$ is locally connected. It is not locally path connected since the only path connected subsets are singletons.
\end{example}

\begin{example}
    The Warsaw circle: connect the topologist's sine curve to the unit circle. This is path connected but not locally connected.
\end{example}

\begin{center}
\begin{tabular}{|l|c|c|}
\hline
\textbf{Ex} & \textbf{Connected} & \textbf{Not connected} \\
\hline
Locally connected
& intervals in $\R$
& $[-1,0)\,\cup\,(0,1]$ \\
\hline
Not locally connected
& topologist's sine curve; $[0,1]\times 0 \cup \Q \times [0,1]$
& $\Q$ \\
\hline
\end{tabular}
\end{center}



\begin{example}
    The broom space: consider the subset of $\R^2$ given by
    \[B=\bigcup_{n\geq 1}\left\{\left(x,\frac{x}{n}\right)|x\in [0,1]\right\}\cup \{(x,0)|x\in[0,1]\}.\]
    Then, $B$ is path connected but not locally path connected at $(1,0)$.
\end{example}

\begin{theorem}
    A space is locally (path) connected iff for any open $U\subset X$, each (path) component $P$ of $U$ is open.
\end{theorem}
\begin{proof}
    Let's prove the path version. 

    ($\Rightarrow$) Let $P$ be a path component of an open set $U\subset X$. For any $x\in P$, there exists a path connected open neighborhood $V$ of $x$ with $V\subseteq U$. Then, $V\subseteq P$ since $P$ is a path component. Thus, $P$ is open.

    ($\Leftarrow$) Given any open neighborhood $U$ of $x\in X$, the path component of $U$ containing $x$ is an open path connected neighborhood of $x$ contained in $U$.
\end{proof}

\begin{theorem}
    If $X$ is locally path connected, then the components and path components of $X$ coincide.
\end{theorem}
\begin{proof}
    Let $P$ be a path component contained in a component $C$. If $P\neq C$, then let $Q$ be the union of all other path components contained in $C$. By the theorem above, we know $C$ is open. Then, by the theorem again, $P$ and $Q$ are open in $C$ and hence separate $C$, contradicting connectedness of $C$.
\end{proof}


\begin{corollary}
    A connected, locally path connected space $X$ is path connected.
\end{corollary}
\begin{proof}
    So $X$ is a single component, which is also a single path component.
\end{proof}

\begin{corollary}
    An open connected subset of $\R^n$ is path connected.
\end{corollary}
\begin{proof}
    Any open subset of $\R^n$ is locally path connected.
\end{proof}


\subsection*{Totally disconnected spaces}

\begin{definition}
    A space $X$ is \emph{totally disconnected} iff the components of $X$ are singletons.
\end{definition}

\begin{example}
    The space $\Q$ of rational numbers (with the subspace topology from $\R$) is totally disconnected.

    Its complement $\R\setminus\Q$ is also totally disconnected.
\end{example}

\begin{example}
    The Cantor set is totally disconnected.
\end{example}

We are actually going to prove this in two ways.

As we know that the Cantor set topology can be generated by an ultrametric, it turns out that the above example is a special case of the following general fact.

\begin{theorem}
    Every ultrametric space is totally disconnected.
\end{theorem}

\begin{proof}
    Let $X$ be an ultrametric space with ultrametric $d$. We need to show that the components of $X$ are singletons. Let $x,y\in X$ be distinct points. Let $r = {u(x,y)} > 0$. Then, $B_r(x)$ is an open neighborhood of $x$. I claim that $X\setminus B_r(x)$ is an open neighborhood of $y$. Indeed, for any point $z\in X\setminus B_r(x)$, we have $u(x,z) \geq r > 0$. Pick $r_0\in(0,r)$. Then, for any $z'\in B_{r_0}(z)$, we have that $u(x,z')=u(x,z)\geq r$. Hence $z'\notin B_r(x)$ and thus $B_{r_0}(z)\subseteq X\setminus B_r(x)$. This shows that $X\setminus B_r(x)$ is open and hence an open neighborhood of $y$. This implies that $x$ and $y$ lie in different components. Since $x,y\in X$ are arbitrary, we conclude that the components of $X$ are singletons.
\end{proof}


\begin{theorem}
    Every product of totally disconnected spaces is totally disconnected.
\end{theorem}
\begin{proof}
    Suppose $C$ is a connected subset of $X=\prod_{\alpha\in A}X_\alpha$ with each $X_\alpha$ totally disconnected. Then, for any $\alpha\in A$, the projection $\pi_\alpha(C)$ is connected in $X_\alpha$ and hence a singleton. Thus, $C$ is a singleton.
\end{proof}

Now, as $\{0,1\}$ is a totally disconnected space, we have that $C\cong \{0,1\}^\N$ is totally disconnected.
\newpage
